\hypertarget{appendix-f414bd4ed05c}{}
\section*{Interpretation contract}
\section{Interpretation contract}\label{interpretation-contract}

\textbf{Curriculum depth:} Developed · \textbf{First seen:} C001 · \textbf{Latest development:} C003

\subsection{Definition and mechanism}\label{definition-and-mechanism}

The width, layout, type, and conventions that turn an otherwise ambiguous bit pattern into data.

For \texttt{0xAD}, unsigned interpretation gives 173; eight-bit two's complement gives −83. A contract must state width, signedness, field positions, byte order, allowed values, and any padding policy. Reinterpretation changes meaning; conversion may change the stored bits.

\subsection{Prerequisites}\label{prerequisites}

\begin{itemize}
\tightlist
\item
  \hyperref[appendix-ddbb69998a2e]{Bit} --- A bit has two possible logical states.
\item
  \hyperref[appendix-280445746493]{Bit significance} --- LSB means least significant bit, normally position 0 with weight 1.
\end{itemize}

\subsection{Used by}\label{used-by}

\begin{itemize}
\tightlist
\item
  \hyperref[appendix-1d0dbe8924cc]{Fixed-point representation} --- uses this mechanism to define its representation, error or validation contract.
\item
  \hyperref[appendix-53a8cfb25f8d]{Offset-binary coding} --- uses this mechanism to define its representation, error or validation contract.
\item
  \hyperref[appendix-cbe594d1ab0a]{Integrity and authenticity} --- uses this mechanism to define its representation, error or validation contract.
\item
  \hyperref[appendix-34ee45f161c1]{Bit casting and representation} --- uses interpretation contract to establish this mechanism's representation or policy.
\item
  \hyperref[appendix-45303d03796a]{Integral promotion} --- uses interpretation contract as part of its representation or reasoning.
\item
  \hyperref[appendix-1f9963e1a5d1]{Bit packing} --- uses interpretation contract as part of its representation or reasoning.
\item
  \hyperref[appendix-f558f5b4d8f2]{Undefined behavior} --- uses interpretation contract as part of its representation or reasoning.
\item
  \hyperref[appendix-91d09ebed979]{Serialization} --- uses interpretation contract as part of its representation or reasoning.
\item
  \hyperref[appendix-942ad41afb30]{Compiler IR and machine instructions} --- uses interpretation contract as part of its representation or reasoning.
\item
  \hyperref[appendix-343615f49ddf]{IEEE-754 floating point} --- uses interpretation contract as part of its representation or reasoning.
\item
  \hyperref[appendix-05db24ccc28a]{Exhaustive finite-state testing} --- uses interpretation contract as part of its representation or reasoning.
\item
  \href{}{Storage duration and object lifetime (05-Engine-I, appendix-0b104bcba06b)} --- depends on this concept; see its prerequisite explanation.
\item
  \href{}{Continuous signals and discrete samples (04-Visual-Computing-II, appendix-1077ea8fc737)}
\item
  \href{}{Color encoding and linear-light operations (04-Visual-Computing-II, appendix-a9f83ad40792)}
\end{itemize}

\subsection{Recall and next depth}\label{recall-and-next-depth}

Explain interpretation contract without the label, work the example above, then state the assumption whose removal would change the result.

\textbf{Pending:} ABI, file formats, instruction encodings, pixels, and protocol schemas.

\subsection{Source and navigation}\label{source-and-navigation}

\href{Continuous\%20Notes/Computing\%20Foundations\%20-\%20From\%20Bits\%20to\%20C++\%20and\%20Systems\%20APIs\#From\%20a\%20signal\%20to\%20an\%20interpretation}{C001 teaching chapter} · \href{Continuous\%20Notes/Computing\%20Foundations\%20-\%20From\%20Bits\%20to\%20C++\%20and\%20Systems\%20APIs\#Arithmetic\%20becomes\%20a\%20policy}{C002 development} · \hyperref[appendix-fd090d982a15]{Part definitions and mechanisms} · \href{}{Dependency Map (02-Foundations-II, appendix-4ad47bdd269a)}

\href{Continuous\%20Notes/Computing\%20Foundations\%20-\%20From\%20Bits\%20to\%20C++\%20and\%20Systems\%20APIs\#Where\%20the\%20bytes\%20live}{C003 development}
